\subsection{深度的同调定义}

设 A 为一个环，J 为 A 的一个理想，M 为一个 A-模。

定义 1. M 关于 J 的深度称为并记作 \(\mathrm{prof}_{\mathrm{A}}(\mathrm{J};\mathrm{M})\) ，或 \(\mathrm{prof}(\mathrm{J};\mathrm{M})\) ，即 \(\mathbf{N}\cup\{+\infty\}\) 中由满足下述条件的整数 \(n\) 组成的集合的下确界：\(\mathrm{Ext}_{\mathrm{A}}^{n}(\mathrm{A}/\mathrm{J},\mathrm{M})\) 非零。

当环 A 为局部环时，把 M 的深度简单地称为并记作 \(\text{prof}_{\mathrm{A}}(\mathrm{M})\) 或 \(\text{prof}(\mathrm{M})\) 的，就是 M 关于 A 的极大理想 \(m_{A}\) 的深度；而局部环 A 的深度则指 A-模 A 的深度。

\(\operatorname{Si} \operatorname{prof}_{\mathrm{A}}(\mathrm{J}; \mathrm{M}) = +\infty\) ，则 \(\operatorname{Ext}_{\mathrm{A}}^{i}(\mathrm{A}/\mathrm{J}, \mathrm{M}) = 0\) 对一切 \(i\) 成立。\(\operatorname{Si} \operatorname{prof}_{\mathrm{A}}(\mathrm{J}; \mathrm{M}) = r < +\infty\) ，则 \(\operatorname{Ext}_{\mathrm{A}}^{i}(\mathrm{A}/\mathrm{J}, \mathrm{M}) = 0\) 对 \(i < r\) 成立，且 \(\operatorname{Ext}_{\mathrm{A}}^{r}(\mathrm{A}/\mathrm{J}, \mathrm{M}) \neq 0\) 。

注.--- 1) 设 A-模 M 有限生成且 M = JM，即 Supp(M) ∩ V(J) = ∅ (II, § 4, 第 4 节, 命题 18 的推论)。此时 prof\_A(J;M) 等于 +∞：事实上，理想 Ann(M) + J 这时等于 A（同上引文），且对一切 i 含于 Ext\_A\^{}i(A/J;M) 的零化子之中。下文（第 3 节, 定理 1 的推论 1）将看到，反之，若理想 J 有限生成，则 prof\_A(J;M) = +∞ 蕴含 M = JM。

\begin{enumerate}
\def\labelenumi{\arabic{enumi})}
\setcounter{enumi}{1}
\tightlist
\item
  为使 \(\mathrm{prof}_{\mathbf{A}}(\mathbf{J};\mathbf{M})\) 为零，必须且只需 \(\mathrm{Hom}_{\mathbf{A}}(\mathbf{A} / \mathbf{J},\mathbf{M})\) 非零，即 M 具有一个被 J 零化的非零元素；特别地，当 \(\mathrm{Ass(M)}\cap \mathrm{V(J)}\neq \emptyset\) 时即属此情形。若环 A 为 Noether 环且 A-模 M 有限生成，则下列条件等价 (IV, § 1, 第 4 节, 命题 8)：
\end{enumerate}

\begin{enumerate}
\def\labelenumi{\alph{enumi})}
\item
  \(\mathrm{prof}_{\mathrm{A}}(\mathrm{J};\mathrm{M}) = 0\)
\item
  对每个 \(x \in J\) ，乘以 x 的映射 \(x_{\mathbb{M}}\) 都不是单射；
\end{enumerate}

\[
\mathrm{c)} \text {   we have   } \operatorname{Ass} (\mathrm{M}) \cap \mathrm{V} (\mathrm{J}) \neq \varnothing .
\]

\begin{enumerate}
\def\labelenumi{\arabic{enumi})}
\setcounter{enumi}{2}
\item
  根据注 2，为使局部环 \(\Lambda\) 的深度为零，必须且只需存在 A 的非零元素 \(x\) 使 \(\mathfrak{m}_{\mathrm{A}}x = 0\) 。若 A 不是域，则元素 \(x\) 一旦满足 \(\mathfrak{m}_{\mathrm{A}}x = 0\) 即不可逆，故属于 \(\mathfrak{m}_{\mathrm{A}}\) ，从而满足 \(x^{2} = 0\) 。因此，维数 \(\geqslant 1\) 的约化局部环的深度 \(\geqslant 1\) 。
\item
  设 \((\mathbf{M}_{\iota})_{\iota \in I}\) 为一族 A-模。根据 A, X, 第 89 页, 命题 7，有 \(\operatorname{prof}(\mathbf{J};\prod_{\iota \in I}\mathbf{M}_{\iota}) = \inf_{\iota \in I}\operatorname{prof}(\mathbf{J};\mathbf{M}_{\iota})\) 。
\end{enumerate}

命题 1.--- 设 A 为环，J 为 A 的理想，且 0 → M' → M → M'' → 0 为 A-模的正合列。令

\[
p ^ {\prime} = \operatorname{prof} (\mathrm{J}; \mathrm{M} ^ {\prime}), p = \operatorname{prof} (\mathrm{J}; \mathrm{M}), p ^ {\prime \prime} = \operatorname{prof} (\mathrm{J}; \mathrm{M} ^ {\prime \prime}).
\]

这时必处于以下三种两两互斥的情形之一：

\[
p ^ {\prime} = p \leqslant p ^ {\prime \prime}, \quad p = p ^ {\prime \prime} <   p ^ {\prime}, \quad p ^ {\prime \prime} = p ^ {\prime} - 1 <   p.
\]

考察与 A/J 及上述正合列相伴的扩张模正合列 (A, X, 第 92 页, 定理 2)。排除 \(p = p' = p'' = +\infty\) 的情形；此时该序列中存在第一个非零模，且紧随其后的模也非零。由此得到以下三种可能：

\begin{enumerate}
\def\labelenumi{\alph{enumi})}
\item
  第一个非零模是 \(\operatorname{Ext}_{\Lambda}^{p'}(A/J, M')\) 。于是 \(p' = p \leqslant p''\) 。
\item
  第一个非零模是 \(\operatorname{Ext}_{\Lambda}^{p}(\mathrm{A} / \mathrm{J},\mathrm{M})\) 。于是 \(p = p'' < p'\) 。
\item
  第一个非零模是 \(\mathrm{Ext}_{\mathrm{A}}^{p''}(\mathrm{A} / \mathrm{J},\mathrm{M}'')\) 。于是 \(p'' + 1 = p' \leqslant p\) 。
\end{enumerate}

注 5.--- 设 \(p = p'\) ，且出现在命题 1 的正合列中的单射 \(u: M' \to M\) 属于 J Hom \(_A(M', M)\) 。于是有 \(a\) \(p'' = p - 1\) 。事实上，该假设蕴含映射 Ext \(_A^i(1_{A/J}, u)\) 对每个整数 \(i\) 均为零；这就排除了上面所考虑的情形 \(a\) )。

命题 2.--- 设 A 为环，J 为 A 的理想，M 为 A-模，N 为被 J 的某个幂所零化的 A-模。则 \(\mathrm{Ext}_{\mathrm{A}}^{i}(\mathrm{N},\mathrm{M}) = 0\) 对每个整数 \(i < \mathrm{prof}_{\mathrm{A}}(\mathrm{J};\mathrm{M})\) 成立。

首先设 JN = 0，并对整数 \(i < \mathrm{prof}_{\mathrm{A}}(\mathrm{J};\mathrm{M})\) 作归纳论证。断言对 \(i < 0\) 显然。把 N 视为 (A/J)-模，并选取 (A/J)-模的一个正合列

\[
0 \rightarrow \mathrm{K} \longrightarrow (\mathrm{A} / \mathrm{J}) ^ {(\mathrm{I})} \longrightarrow \mathrm{N} \rightarrow 0.
\]

由此导出扩张模的一个正合列

\[
\operatorname{Ext} _ {\mathrm{A}} ^ {i - 1} (\mathrm{K}, \mathrm{M}) \longrightarrow \operatorname{Ext} _ {\mathrm{A}} ^ {i} (\mathrm{N}, \mathrm{M}) \longrightarrow \operatorname{Ext} _ {\mathrm{A}} ^ {i} ((\mathrm{A} / \mathrm{J}) ^ {(1)}, \mathrm{M}).
\]

A-模 \(\mathrm{Ext}_{\mathrm{A}}^{i-1}(\mathrm{K},\mathrm{M})\) 由归纳假设为零，而 A-模 \(\mathrm{Ext}_{\mathrm{A}}^{i}((\mathrm{A}/\mathrm{J})^{(1)},\mathrm{M})\) 同构于 \(\mathrm{Ext}_{\mathrm{A}}^{i}(\mathrm{A}/\mathrm{J},\mathrm{M})^{\mathrm{I}}\) (A, X, 第 89 页, 命题 7)，后者依深度的定义而为零。于是有 \(\mathrm{Ext}_{\mathrm{A}}^{i}(\mathrm{N},\mathrm{M}) = 0\) 。

转入一般情形，并对最小的整数 m \textgreater{} 0 作归纳论证，其中 m 满足 \(J^{m}N = 0\) 。我们刚才已处理 m = 1 的情形。设 m \textgreater{} 1 ，并设整数 \(i < \text{prof}_{\text{A}}(\text{J}; \text{M})\) 。考察正合列

\[
\operatorname{Ext} _ {\Lambda} ^ {i} (\mathrm{N/JN,M}) \longrightarrow \operatorname{Ext} _ {\Lambda} ^ {i} (\mathrm{N,M}) \longrightarrow \operatorname{Ext} _ {\Lambda} ^ {i} (\mathrm{JN,M})
\]

它由正合列 \(0 \rightarrow JN \rightarrow N \rightarrow N/JN \rightarrow 0\) 导出。两端的模由归纳假设为零，因为 N/JN 与 JN 都被 \(J^{m-1}\) 零化。于是有 \(\operatorname{Ext}_{\Lambda}^{i}(N,M) = 0\) ，此即所要证明的。

推论 1.--- 设 m 为整数 \textgreater0 ，J' 为 A 的包含 J\^{}m 的理想。则 \(\operatorname{prof}_{\mathrm{A}}(\mathrm{J};\mathrm{M}) \leqslant \operatorname{prof}_{\mathrm{A}}(\mathrm{J}';\mathrm{M})\) 。

事实上，\(J^{m}\) 零化 A-模 \(A/J'\) ，故 \(\operatorname{Ext}_{A}^{i}(A/J', M)\) 对每个整数 \(i < \operatorname{prof}_{A}(J; M)\) 均为零（命题 2）。

推论 2.--- 设理想 J 有限生成，并设 J' 为 A 的理想且满足 V(J) ⊃ V(J')。

\begin{enumerate}
\def\labelenumi{\alph{enumi})}
\item
  有 \(a \operatorname{prof}_{\mathrm{A}}(\mathrm{J}; \mathrm{M}) \leqslant \operatorname{prof}_{\mathrm{A}}(\mathrm{J}'; \mathrm{M})\) 。
\item
  若理想 J' 有限生成且 V(J) = V(J')，则有 \(\text{prof}_{A}(J; M) = \text{prof}_{A}(J'; M)\) 。
\end{enumerate}

根据 II, § 4, 第 3 节, 命题 11 的推论 2 及 § 2, 第 6 节, 命题 15，存在整数 m \textgreater{} 0 使 \(J^{m} \subset J'\) 。于是断言 a) 由推论 1 得出，而断言 b) 又由它导出。

当理想 J 不是有限生成时，推论 2 可能不成立（习题 2）。

